% Speaker script for slides_defense.tex (colloquium, Wed 16 Dec 2026, 13:00), English + Japanese.
% Build: lualatex speaker_notes_defense.tex (twice), or
%        .\build_slides.ps1 -Tex speaker_notes_defense.tex -MaxPages 30 on IKMHIWI03.
% EN is what I say; JA is the same in Japanese, for practice, plus hints in grey.
% Numbers are those of the thesis; check them again if the thesis changes.
\documentclass[11pt,a4paper]{article}
\usepackage{luatexja}
\usepackage[margin=22mm]{geometry}
\usepackage{amsmath,amssymb}
\usepackage{xcolor}
\usepackage{enumitem}
\definecolor{c3}{RGB}{30,80,160}
\setlength{\parindent}{0pt}
\setlength{\parskip}{4pt}
\newcounter{sl}
\newcommand{\slide}[2]{\stepcounter{sl}\vspace{8pt}%
  {\color{c3}\large\bfseries Slide \thesl\ --- #1}\hfill{\small #2}\par\vspace{2pt}}
\newcommand{\EN}{\textbf{EN}\quad}
\newcommand{\JA}{\par\textbf{JA}\quad}
\newcommand{\hint}[1]{\par{\color{gray}\small ヒント：#1}}

\begin{document}
{\Large\bfseries Speaker script --- master's thesis colloquium, 16 December 2026}\\[2pt]
Prof.\ Junker, Dr.\ Wangenheim, Dr.\ Soleimani, Dr.\ Geisler.
16 slides plus a 3-slide appendix, about 20 minutes. The times are the
running total at the end of each slide.

% ------------------------------------------------------------------
\slide{Title}{0:30}
\EN Good afternoon. Thank you for being here. My thesis is about coupling
biofilm growth into an existing finite element framework: bringing a
multi-species point model into the biofilm element of the institute, so that
the composition, the growth and the stresses of a biofilm can be computed
together.
\JA こんにちは。お越しいただきありがとうございます。私の修論は、多菌種の点モデルを
研究所のバイオフィルム要素に組み込み、組成・成長・応力を一緒に計算できるようにする、
というものです。

\slide{Motivation}{1:30}
\EN Peri-implantitis affects about one in five implant patients. It is driven
by multi-species biofilms, and a shift towards pathobionts such as
\textit{Porphyromonas gingivalis} goes with disease. A biofilm is also a
mechanical body: where it grows unevenly, or against a stiff substrate, it
carries stresses. Two kinds of model exist. A point model describes the
composition, but has no space. A growth field describes space and stress, but
does not distinguish species. My question was whether both can be computed
together, inside a framework that already works.
\JA インプラント周囲炎は、インプラント患者の約5人に1人に起きます。原因は多菌種の
バイオフィルムで、\textit{P.~gingivalis} のような病原性の菌が増えると悪化します。
バイオフィルムは力学的な物体でもあり、成長が不均一だったり硬い面に押し付けられたり
すると応力を持ちます。モデルは2種類あります。点モデルは組成を表すが空間がない。
成長場は空間と応力を表すが菌種を区別しない。この2つを、すでに動いている枠組みの中で
一緒に計算できるか、が問いです。

\slide{Overview}{2:30}
\EN This is the whole work on one slide. On the left, the growth field of
Klempt et al.\ 2024 in the IKM element gives the amount of biofilm at every
Gauss point. In the middle, one copy of the point model of Klempt et al.\
2026 sits at every Gauss point and decides only the composition. On the
right, growth and stress use the amount only. The coupling is one-way: the
composition is an output, it does not act back. The talk follows the three
parts: calibrating the point model, the growth law in ANSYS, and the
composition at every Gauss point.
\JA 全体を1枚にした図です。左：Klempt 2024 の成長場が、各 Gauss 点のバイオフィルムの
量を決めます。中央：Klempt 2026 の点モデルを各 Gauss 点に1つずつ置き、組成だけを決めます。
右：成長と応力は量だけを使います。一方向の結合で、組成は出力であって戻りません。
発表はこの順に3部です。
\hint{最初の「何をしたか」を30秒で言えるようにしておく。}

\slide{Both models from the extended Hamilton principle}{4:00}
\EN Both models come from the extended Hamilton principle. The growth field
has the biofilm variable $\phi$, the nutrient $c$ and the growth $\alpha$,
with $\mathbf F=\mathbf F_e\mathbf F_g$ and $\mathbf F_g=\alpha\mathbf I$;
stress comes only from the elastic part. For the point model there is no
displacement, so the principle reduces to its local form, the principle of
the minimum of the dissipation potential: free energy rate plus dissipation
plus constraint, stationary with respect to the rates and the multiplier. One
consequence is that the interaction matrix $\mathbf A$ is symmetric.
\JA どちらのモデルも拡張 Hamilton 原理から出ます。成長場は $\phi$、$c$、$\alpha$ を持ち、
$\mathbf F_g=\alpha\mathbf I$、応力は弾性部分からだけ出ます。点モデルには変位がないので、
原理はその局所形、最小散逸ポテンシャルの原理になります。結果として相互作用行列
$\mathbf A$ は対称です。
\hint{Junker 先生はここを聞く。$\Psi$, $\Delta$, $C$ を板書できるように。}

\slide{Part 1: calibrating the point model}{5:30}
\EN The first part calibrates the point model. Heine et al.\ cultured five
species under four conditions, commensal or dysbiotic, static or in a flow
reactor, and measured the composition from day 1 to day 21. I estimated the
15 interaction coefficients for each condition with transitional Markov chain
Monte Carlo, with ten thousand samples. The forward model runs in JAX on the
GPU, about 150 times faster than on the CPU, which makes this practical. The
lines here are the fits against the measured points.
\JA 第1部は点モデルの較正です。Heine らの5菌種・4条件のデータを使い、各条件で15個の
相互作用係数を TMCMC で推定しました。サンプルは1万個です。順モデルを JAX で GPU に
載せたので、CPU より約150倍速くなり、現実的な時間で回せます。

\slide{Part 1: what the calibration gives}{7:00}
\EN The calibration gives one interaction matrix per condition. Under
dysbiotic flow, the link from \textit{F.~nucleatum} to \textit{P.~gingivalis}
becomes strongly cooperative. On the right, the pH, which I did not use in
the calibration, is predicted with $R^2=0.78$. This leads to a prediction
that can be tested: without \textit{F.~nucleatum}, the late growth of
\textit{P.~gingivalis} under dysbiotic flow should be suppressed.
\JA 条件ごとに相互作用行列が得られます。dysbiotic の流れ条件では、\textit{F.~nucleatum}
から \textit{P.~gingivalis} への協力が強く出ます。右は pH で、較正には使っていませんが、
$R^2=0.78$ で予測できました。ここから検証できる予測が出ます：\textit{F.~nucleatum} を
除けば、終盤の \textit{P.~gingivalis} の増加は抑えられるはずです。

\slide{Part 2: the growth law in ANSYS}{8:30}
\EN The second part brings the growth law into ANSYS. I implemented
$\mathbf F_g=\alpha\mathbf I$ as a user material, with the stiffness of
Klempt et al.\ 2024: ten pascal, Poisson's ratio 0.49, weighted by $\phi^2$.
The small floor $f$ for the void is my addition. On one element I compared
free and constrained growth with closed forms: free growth gives zero
stress, and constrained growth gives the predicted pressure.
\JA 第2部は成長則を ANSYS に入れた部分です。$\mathbf F_g=\alpha\mathbf I$ をユーザー
材料として実装し、剛性は Klempt 2024 の値（10 Pa、0.49、$\phi^2$ で重み付け）にしました。
空隙の下限 $f$ は私が足したものです。1要素で、自由な成長と拘束した成長を閉じた形の解と
比べ、自由なら応力ゼロ、拘束なら予測どおりの圧力になりました。
\hint{$f$ は論文にないと自分から言う。}

\slide{Part 2: coupling into the IKM element}{9:30}
\EN This is the example model: a two-millimetre cube with a seed of biofilm
in the middle, refined up to $24^3$ elements. The element's own field for
$\phi$ is unchanged. What I added is the growth law at every Gauss point:
$\phi$ gives $\alpha$, $\alpha$ gives $\mathbf F_g$, and that gives the stress.
The diffusion coefficient $\beta$ is Table~2 of the paper converted to this
cube; this conversion is an assumption, and I report its sensitivity.
\JA 例題は 2 mm の立方体で、中央にシードがあります。$24^3$ まで細かくしました。
要素の $\phi$ の場はそのままで、私が足したのは各 Gauss 点の成長則です。
$\beta$ は論文の Table~2 をこの立方体に換算した値で、仮定なので感度も示しています。

\slide{Part 2: the coupled run against exact solutions}{10:30}
\EN Where $\phi$ has no gradient, the equations have an exact solution:
$\phi$ grows as a hyperbolic sine and $\alpha-1$ as a hyperbolic cosine minus
one. The explicit update has its own exact solution, and ANSYS follows it to
$3\times10^{-11}$ in $\alpha$. So the coupling inside ANSYS does what the
equations say.
\JA 勾配のないところでは厳密解があり、$\phi$ は sinh、$\alpha-1$ は cosh $-1$ になります。
陽解法の更新にはそれ自身の厳密解があり、ANSYS は $\alpha$ で $3\times10^{-11}$ まで一致します。

\slide{Part 2: stresses with the published stiffness}{12:00}
\EN With the published stiffness, the seed is compressed, about
$-2.5\times10^{-4}$ pascal mean stress, and the first layer of elements around
it is in tension. This is the ring of tension that Klempt et al.\ report.
The stresses are small because the stiffness is ten pascal, and they scale
with it. What carries over is the pattern, not the magnitude.
\JA 論文の剛性では、シードは圧縮（平均 $-2.5\times10^{-4}$ Pa）、周りの1層目は引張です。
Klempt らが報告している引張の輪と同じです。応力が小さいのは剛性が 10 Pa だからで、
剛性に比例します。意味があるのは大きさではなく分布です。
\hint{Wangenheim 先生は「この大きさに意味はあるか」を聞きそう。}

\slide{Part 2: how far the numbers can be trusted}{13:30}
\EN How far can these numbers be trusted? In time, the update is first
order; at my time step the error is about one percent. In space, the seed
stress converges at an observed order of about two, with an error of about
two percent on $24^3$ elements, estimated from three meshes. The growth,
compared at fixed positions, changes by less than four percent from $16^3$ to
$24^3$. And the explicit field update is stable only up to $\beta\Delta
t/h^2$ of about 0.15 in three dimensions, which is stricter than the
condition given in the literature. All runs in the thesis are below it.
\JA 時間は1次で、今の刻みの誤差は約 1 \%。メッシュは約2次で収束し、$24^3$ の誤差は約
2 \%（3つのメッシュからの推定）。$\alpha$ は位置を固定して比べると $16^3\to24^3$ で 4 \% 以内。
陽解法の安定の境目は 3 次元で約 0.15 で、文献の条件より厳しい。修論の実行はすべてその下。

\slide{Part 3: composition at every Gauss point}{15:00}
\EN The third part puts the composition at every Gauss point, for two species
and the cases of Klempt et al.\ 2026. Every call from ANSYS to the point model
is reproduced in Python bit for bit. In the seed, case three keeps
coexistence, with a share of 0.67, and in case six one species wins, with a
share of 0.02. The time link $s$ and the cap $\phi_{\rm cap}$ are not in the
papers; they are my assumptions, each with its sensitivity.
\JA 第3部では、2菌種で各 Gauss 点に組成を持たせました。ANSYS から点モデルへの呼び出しは
Python とビット単位で一致します。シードでは case 3 は共存（0.67）、case 6 は一方が勝つ
（0.02）で、結果の型は保たれます。$s$ と $\phi_{\rm cap}$ は論文にない仮定です。

\slide{Part 3: composition varies in space with the local nutrient}{16:30}
\EN If the nutrient is held at one face and consumed in the seed, the
composition varies in space even without spreading. The strength is the
Thiele number. In case six, at $\Lambda=4$, the share is 0.03 near the
nutrient and 0.44 inside the seed; case three hardly changes. This is
something the element cannot show with a constant nutrient level.
\JA 栄養を一つの面に与えてシードで消費させると、広がらなくても組成が空間で変わります。
強さは Thiele 数です。case 6、$\Lambda=4$ では、栄養側で 0.03、シードの奥で 0.44。
case 3 はほとんど変わりません。

\slide{Limitations}{18:00}
\EN The main limitations. This is verification, not validation: no measured
stress of a growing oral biofilm was available. The coupling is one-way. The
front term of the element is inactive, so there is no spreading towards the
nutrient in ANSYS. Three values are assumptions. The calibrated five species
and the element have not yet met in one computation. And a symmetric matrix
cannot describe directed interactions.
\JA 主な限界です。検証はしたが妥当性確認はしていない（応力の測定値がない）。結合は一方向。
ANSYS では前線の項が働かず、栄養へ向かう広がりはない。3つの値は仮定。較正した5菌種と
要素はまだ一緒になっていない。対称な行列は向きのある相互作用を表せない。
\hint{限界はこちらから先に言う。聞かれる前に言えば強い。}

\slide{Conclusion and outlook}{19:30}
\EN To conclude: I calibrated a point model on in-vitro data, and it predicts
the pH out of sample; I brought the growth law of Klempt et al.\ 2024 into
the IKM element and verified it; and I put the composition at every Gauss
point for two species, keeping the published outcomes. At Keio I will
continue with the five calibrated species in the element, the two-way
coupling, a working front term, and the uncertainty carried to the stress.
Thank you.
\JA まとめ：点モデルを較正し pH を予測できた。Klempt 2024 の成長則を IKM 要素に入れて
検証した。2菌種で各 Gauss 点に組成を持たせ、論文の結果を保った。慶應では、5菌種を要素に、
双方向の結合、前線の項、不確かさを応力まで、を続けます。ありがとうございました。

\slide{References}{---}
\EN (Only if asked. Leave this slide up during the questions.)
\JA （聞かれたときだけ。質疑の間はこのページを出しておく。）

% ------------------------------------------------------------------
\vspace{10pt}
{\large\bfseries Appendix slides (for questions)}\par
\begin{itemize}[leftmargin=*]
  \item \textbf{A, notation}: 記号を聞かれたとき。
  \item \textbf{B, growth field}: 場の式・材料・消費項 $g\phi c$ を聞かれたとき。
  \item \textbf{C, point model and one substep}: 点モデルの式と、1サブステップの流れ。
\end{itemize}

\end{document}
